\scanpage{149}
% archive sheet number: 148
``de type $\mathbb Q_\ell$'' si il est isomorphe \unsure{à|:} un \deleted{\unsure{limite}} système inductif $(G_n)_{n\in\mathbb N}$, où pour tout $n$, on a $G_n=G_0$ $G$ étant ``de type $\deleted{\unsure{\mathbb Q}}\mathbb Z_\ell$'', et $\unsure{\varphi_{nm}|\varphi_{mn}}$ (pour $m\geq n$) étant la multiplication par $\ell^{m-n}$. Si $S$ est connexe, de pt géométrique $\xi$, alors la catégorie \deleted{de} \unsure{ces|des} $G$ est équivalente à la catégorie des espaces vectoriels \inserted{$\unsure{N|V}$} de dim finie sur $\mathbb Q_\ell$, munis d'un hom continu de groupes topologiques $\pi_1(S,\xi)\longrightarrow\deleted{\unsure{\operatorname{GL}(N)}}\operatorname{Aut}(N)$. [Noter que $\pi_1(S,\xi)$ étant compact, il y \unsure{aura|a} un ``réseau'' invariant ...]. Soit de façon générale \inserted{$\unsure{\mathcal V|\mathscr V}_\ell(S)$ la catégorie des objets de $\operatorname{Ind}\operatorname{Pro}\unsure{\mathscr G|\mathcal G}(S)$ qui sont ``de type $\mathbb Q_\ell$'', $\unsure{\mathcal M|\mathscr M}_\ell(S)$ la catégorie des objets} \deleted{\unsure{la catégorie des} \illegible{1}} de $\operatorname{Pro}\mathscr G(S)$ qui sont ``de type $\mathbb Z_\ell$'', alors \deleted{\unsure{on}}
\[
M\unsure{\rightsquigarrow|\longrightarrow|\longmapsto}M\otimes_{\mathbb Z_\ell}\mathbb Q_\ell=(M_n=M,\;M_m\xrightarrow{\ell^{m-n}}M_n)
\]
est un foncteur
\[
\mathcal V_\ell(S)\longrightarrow\mathcal M_\ell(S)
\]
qui est essentiellement surjectif, \inserted{et \uline{fidèle}}, (pas pleinement fidèle en général). Les
